% Kolmo: DCC-style paper draft. Build: make -C docs/paper   (see docs/paper/README.md)
% DRAFT. Numbers come from docs/paper/generated/*.tex, which figures/make_figures.py writes from docs/paper/data/.
% \drafttrue prints a source tag after each number in the running text and shows TODO boxes; \draftfalse hides both.
\documentclass[conference]{IEEEtran}
\usepackage{iftex}
\ifXeTeX
  \usepackage{fontspec}
  \setmainfont{texgyretermes}[Extension=.otf, UprightFont=*-regular, BoldFont=*-bold, ItalicFont=*-italic, BoldItalicFont=*-bolditalic]
\else
  \usepackage[T1]{fontenc}
  \usepackage{newtxtext,newtxmath}
\fi
\usepackage{amsmath}
\usepackage{booktabs}
\usepackage{graphicx}
\usepackage{xcolor}
\usepackage{url}
\usepackage{cite}
\usepackage{pgfplots}
\pgfplotsset{compat=1.18}
\usepgfplotslibrary{groupplots}

\newif\ifdraft
\drafttrue
\ifdraft
  \newcommand{\src}[1]{\textsuperscript{\tiny\color{gray}\texttt{[#1]}}}
  \newcommand{\todo}[1]{\noindent\fcolorbox{red}{red!6}{\parbox{\dimexpr\linewidth-2\fboxsep-2\fboxrule}{\footnotesize\textbf{TODO:} #1}}}
  \newcommand{\todoin}[1]{{\color{red}\footnotesize[TODO: #1]}}
\else
  \newcommand{\src}[1]{}
  \newcommand{\todo}[1]{}
  \newcommand{\todoin}[1]{}
\fi

\input{generated/numbers.tex}

% ---- plot styles (figures/make_figures.py writes the coordinates)
\definecolor{ckolmo}{HTML}{C0392B}
\definecolor{czstd}{HTML}{1F5FA8}
\definecolor{cxz}{HTML}{2E8B57}
\definecolor{cbrotli}{HTML}{8E44AD}
\definecolor{cother}{HTML}{555555}
\pgfplotsset{
  paretoaxis/.style={width=\columnwidth, height=0.62\columnwidth, ymin=1.9, ymax=6.2, grid=major,
    grid style={gray!20}, tick label style={font=\scriptsize}, label style={font=\footnotesize},
    legend style={font=\scriptsize, draw=none, legend columns=4, /tikz/every even column/.append style={column sep=4pt}}},
  frontier/.style={gray!60, dashed, thick, mark=none},
  ppkolmo/.style={ckolmo, very thick, mark=*, mark size=1.8pt},
  ppzstd/.style={czstd, thick, mark=square*, mark size=1.4pt},
  ppbrotli/.style={cbrotli, thick, mark=diamond*, mark size=1.8pt},
  ppxz/.style={cxz, thick, mark=triangle*, mark size=2pt},
  ppkanzi/.style={orange!80!black, thick, mark=pentagon*, mark size=1.6pt},
  pplzfour/.style={cother, mark=o, mark size=1.6pt},
  pplzav/.style={teal, mark=x, mark size=2.2pt, thick},
  ppdeflate/.style={brown, mark=+, mark size=2.2pt, thick},
  ppbwt/.style={black, mark=star, mark size=2.2pt},
  ppcm/.style={magenta!70!black, mark=asterisk, mark size=2.2pt},
}
\tikzset{kolmolabel/.style={font=\tiny\bfseries, text=ckolmo, anchor=south east, inner sep=1pt}}

\begin{document}

\title{Kolmo: Zero-Configuration, Structure-Aware\\Block Routing in a Stable Container}

\author{\IEEEauthorblockN{Adnane Aqartit}
\IEEEauthorblockA{Independent \todoin{affiliation and contact}}}

\maketitle

\begin{abstract}
General-purpose lossless compressors treat their input as a byte stream; format-aware systems such as OpenZL obtain
much smaller output on structured data but need the user to name a format profile or train a compression plan. We
describe Kolmo, a compressor that sits between the two. For every block it detects structure from the bytes alone
(delimited tables, logs, JSON, integer and floating-point arrays, images, executables, genomic reads), chooses a
reversible transform or a small typed transform graph by trial coding on samples, and codes the result with one of
five coder families selected by a single compression level. All choices are recorded per block in a container whose
version-12 layout is specified and frozen, whose blocks decode independently, and whose header declares the memory a
decoder needs before it allocates anything. On one thread of an Apple M5, Kolmo's level~6 is
\numSilKSixR{} on Silesia against \numSilxzNineR{} for \texttt{xz -9}, and decodes \numSilSixVsXzDec{} times faster;
on a panel of \numPanelFiles{} structured files it matches or beats the best untrained OpenZL profile on every file at
its strongest level, without being told the format. We also report where the approach does not help: small inputs,
plain text at fast levels, model weights, decode speed relative to zstd, and memory use.
\todoin{update with x86-64 rerun, trained-OpenZL rematch and tuned levels}
\end{abstract}

\ifdraft
\todo{Draft of 2026-11-23 for DCC 2028. Every number in the running text is a macro from
\texttt{generated/numbers.tex}; a grey tag gives its data file in \texttt{docs/paper/data/}. Three measurement lanes
still running will replace numbers: (1) the x86-64 rerun of all lzbench tables (\S\ref{sec:setup}); (2) the
OpenZL rematch with \texttt{zli train} on public held-out data (\S\ref{sec:openzl}); (3) the encoder tuning of
levels 1--3 and 5--7 (\S\ref{sec:levels}). Nothing here has been submitted.}
\fi

% ======================================================================
\section{Introduction}

Most data that is stored at scale is not prose. It is CSV and JSON exports, server logs, columns of integers and
timestamps, scientific arrays, images and executables. General-purpose compressors---zstd~\cite{zstd_rfc8878},
xz/LZMA~\cite{lzma_xz,lzma_sdk}, brotli~\cite{brotli_rfc7932}, the block-sorting coders~\cite{bwt,bsc,bzip3}---model such
data as a sequence of bytes. They capture repeats and local byte statistics, but not the fact that the fourth column
of a table is a monotone timestamp or that every fourth byte of an array is a floating-point exponent. Domain systems
capture exactly that: columnar formats with their encodings~\cite{parquet,alp,btrblocks,gorilla}, genomic formats
such as CRAM~\cite{cram,cram31}, coders for neural-network weights~\cite{zipnn,dfloat11}. Each works for one kind of
data and must be chosen and configured by whoever stores it.

OpenZL~\cite{openzl,openzl2026,openzl_web} recently made the case for a middle ground: a general framework in which a compressor
is a graph of typed transforms and codecs, described by a self-describing frame so that one universal decoder reads
every configuration. Its best results come from a format profile chosen by the user, or a plan trained on samples of
the data.

This paper describes Kolmo, a compressor built on a narrower premise: \emph{that a useful part of that gain can be had
with no configuration at all}, if the encoder detects structure per block from the bytes, tries a small set of
candidate transforms on samples, and records the outcome in a fixed, specified container. Kolmo is a single-maintainer
open-source project in C (Apache-2.0); version 0.0.2 was released on 2026-11-23.

\paragraph*{Contributions} We do not claim new transforms or a new entropy coder; most of Kolmo's building blocks are
known, and we cite their origins. The contributions are:
\begin{enumerate}
\item A per-block \emph{routing} design that maps detected structure to transforms and coders without user input,
using only the content (never file names), with trial coding on samples as the arbiter (\S\ref{sec:detect}).
\item A container that carries the routing decision per block, with independent blocks, a declared decode working
set, and a written compatibility promise backed by a frozen corpus of frames (\S\ref{sec:container}).
\item A determinism discipline: the same input and level give the same bytes on every CPU, thread count and memory
budget, which we test in continuous integration (\S\ref{sec:determinism}).
\item A single-thread evaluation against 50-odd general-purpose codec settings on the standard corpora and against
OpenZL 0.3.0's format profiles on \numPanelFiles{} structured files, including the cases where Kolmo loses
(\S\ref{sec:eval}, \S\ref{sec:limits}).
\end{enumerate}

% ======================================================================
\section{Related work}\label{sec:related}

\paragraph*{General-purpose coders} LZ77~\cite{lz77} with entropy coding is the basis of deflate, zstd and brotli.
Zstandard~\cite{zstd_rfc8878,zstd_web} pairs an LZ parse with Huffman-coded literals and finite-state-entropy coded
sequences, a form of asymmetric numeral systems~\cite{ans}; its decoder runs at memory speed and is the default
choice of most storage systems. Brotli~\cite{brotli_rfc7932,brotli_tois} adds context modelling of literals and a
static dictionary. LZMA, as used by xz~\cite{lzma_xz} and 7-Zip~\cite{lzma_sdk}, codes an optimal parse with an adaptive
binary range coder; it compresses well and decodes slowly. Block sorting~\cite{bwt}, with modern suffix
sorting~\cite{sais,libsais} and coders such as libbsc~\cite{bsc} and bzip3~\cite{bzip3}, sits between LZMA and
context mixing on text. Kolmo reimplements an LZ4-class coder, a zstd-class coder and an LZMA-class coder from their
published descriptions and ports libbsc's QLFC and LZP stages and libsais to C (\S\ref{sec:coders}).

\paragraph*{Context mixing} The PAQ family~\cite{paq,dce_book} predicts bits by mixing many context models with
online-trained neural networks; lpaq~\cite{lpaq}, zpaq~\cite{zpaq}, cmix~\cite{cmix} and NNCP~\cite{nncp,nncp_v2} explore
the speed and ratio trade-off from practical to extreme, and lead the Large Text Compression
Benchmark~\cite{ltcb}. Kolmo's ultra tier is an lpaq-style integer mixer; it is far from the top of that table and
is not a contribution of this paper.

\paragraph*{Format-aware compression} Columnar formats apply per-column encodings (dictionary, delta, run length,
bit packing) before a general codec~\cite{parquet}; BtrBlocks~\cite{btrblocks} chooses among them per block by
sampling, which is the closest precedent for Kolmo's trial-based selection. Gorilla~\cite{gorilla} and
ALP~\cite{alp} target time series and floating-point columns. CRAM~\cite{cram,cram31} compresses sequencing reads
against a reference; ZipNN~\cite{zipnn} and DFloat11~\cite{dfloat11} exploit the skewed exponents of model weights.
Array containers such as HDF5~\cite{hdf5}, Zarr~\cite{zarr} and Blosc~\cite{blosc} expose byte shuffling and
pluggable codecs, again configured by the user.

\paragraph*{OpenZL} OpenZL~\cite{openzl} is the most direct comparison. It models compression as a directed graph of
codecs over typed streams, serialises the graph into the frame so a single decoder reads any configuration, and
offers format profiles (CSV, little-endian integer arrays, NumPy, pickle, SAO, Parquet, model formats) plus a trainer
that searches for a graph on sample data. Kolmo's graph blocks (\S\ref{sec:graph}) share the idea of a small program
of typed nodes whose inverses the decoder runs, but differ in three ways: the graph is chosen by the encoder per block
without a profile; the set of node types is closed and versioned with the container rather than open; and the
container makes promises about stability, determinism and decode memory that we describe below. OpenZL's design is
more general; Kolmo's is narrower and requires nothing from the user.

% ======================================================================
\section{Container}\label{sec:container}

A Kolmo frame is a 32-byte header, an index of 16-byte block entries, the block payloads, and optional extensions.
The header holds the original length, the block count, a nominal block size, a flag byte, and a one-byte
declaration of the decode working set (\S\ref{sec:memd}); a CRC-32C covers the header and the whole index. Each
index entry gives a block's decoded and stored lengths, a 32-bit checksum, a codec identifier and a transform flag.
Blocks are independent: a decoder may decode them in any order, in parallel, or only those overlapping a requested
byte range.

Every payload starts with a four-byte \emph{plan prefix} $(t,a,b,c)$: a plan type and three parameters (for example,
the delimiter and column count of a table, or the element width and channel count of a byte-plane shuffle). Twelve
plan types are defined (Table~\ref{tab:plans}). The transform flag in the index says which transform actually
produced the stored bytes; it may be zero when the planner chose a plan but the trial found no gain. A decoder reads
the plan and any transform metadata, decodes the coded stream with the block's codec, undoes the transform, and
checks the block checksum. Every length the format carries is bounded before allocation (for example, a transformed
block may exceed its raw length by at most a fixed margin that depends on the transform).

\begin{table}[t]
\caption{Plan types of container version 12 (FORMAT.md \S6.1).}\label{tab:plans}
\centering\footnotesize
\begin{tabular}{@{}rl@{}}
\toprule
$t$ & Plan \\ \midrule
0 & none \\
1 & typed transpose of delimited text \\
2 & byte-plane shuffle (width 2/4/8, channels, delta) \\
3 & typed table with inter-column formulas \\
4 & log templates (word, JSON, minified-JSON tokenisers) \\
5 & 16-bit image predictor \\
6 & fixed-size records \\
7 & x86 call/jump filter \\
8 & word dictionary \\
9 & transform graph (codec 12, \S\ref{sec:graph}) \\
10 & line classes \\
11 & ARM64 branch filter with page phase \\
\bottomrule
\end{tabular}
\end{table}

\paragraph*{Specification and stability} The container layer of version 12 is specified in a normative document
(\texttt{FORMAT.md}) from which an independent Python implementation was written; it checks every frame of a frozen
corpus in continuous integration on Linux x86-64, Linux ARM64 and macOS. The project's compatibility promise is that
every later release decodes every version-12 frame to the same bytes, that identifiers never change meaning, and that
a new feature uses an identifier older decoders refuse, so that an old decoder fails rather than returning wrong
data. The corpus holds 341 frames written by the 0.0.1 and 0.0.2 release binaries
\todoin{count after the format-prose PR merges; was 330 at PR~\#18}. Coder payloads are specified in part; for the
rest the 0.0.2 source is normative until their sections are written.

\paragraph*{Integrity} Block checksums are the low 32 bits of a 64-bit hash written for Kolmo and specified in full
in \texttt{FORMAT.md}. A decoder never returns the bytes of a block whose checksum fails. The decoder has been fuzzed
with libFuzzer \todoin{cite the CI fuzz workflow and OSS-Fuzz status once PR~\#19 lands}.

\subsection{Decode working set}\label{sec:memd}
Many formats declare a window or dictionary size so that a reader can refuse a stream it cannot afford before
allocating (zstd's window descriptor, xz's dictionary size). Kolmo generalises this to all of its coders. The
working set of block $i$ is a function of its index entry and the head of its payload only: the decoded block,
the transform buffer and small state ($2r + r/8 + 16\,\text{KiB} + 1\,\text{MiB}$ for raw length $r$), plus a
codec-specific term (none for the LZ and Huffman coders, the dictionary for the LZMA-class coder, $5t$ for block
sorting, exactly the hash tables named in the stream header for context mixing). The header byte \texttt{memd}
encodes the largest figure in the frame. A reader rejects a frame whose blocks exceed their declaration, compares
the declaration with its memory limit before decoding, and refuses context-mixing tables larger than declared before
allocating them. The declaration is computed from the payloads, not from the free memory of the encoder.

% ======================================================================
\section{Detection and planning}\label{sec:detect}

\paragraph*{Principle} Every decision comes from the bytes; file names and extensions are never read, so identical
bytes always give identical frames. Structure stated in a header (image geometry, executable sections, tar members,
NumPy dtype and shape, safetensors and GGUF tensor tables) is used as a \emph{hint} that moves block boundaries or
fixes a parameter, never as a dependency: each block carries its full plan and the decoder reads no metadata. Header
parsers treat their input as untrusted, bound every offset and count, and treat an unparsable header as absent; they
are tested on thousands of mutated inputs under sanitizers.

\paragraph*{Segmentation} Before planning, the input is cut into parts that deserve separate plans: tar members
(those of at least 64\,KiB with a recognised header, or at least 1\,MiB otherwise), stored ZIP and NPZ members, the
code and data sections of ELF, Mach-O and PE files, runs of one dtype in tensor files, and the sample chunks of WAV
and AIFF audio. Already-compressed content (gzip, xz, zstd, PNG, JPEG, media containers) is recognised by its magic
number and tested with a cheap entropy and LZ probe, and stored blocks are emitted without running the slow coders.

\paragraph*{Text} A part whose first 64\,KiB are at least 95\% printable is planned as text. The planner tests, in
order: very wide delimited tables (more than 256 columns), which go to a transform graph; minified JSON; JSON lines and
pretty-printed JSON; delimited tables of 2--256 fields with at least 95\% of lines alike (comma, tab, semicolon, bar
or space); log lines, by a template trial; and, at higher levels, a word dictionary. Each candidate is kept only if a
trial coding of a sample with a fast coder is smaller than the plain coding by a fixed margin (3\% for tables, 8\% for
JSON and log templates, 1\% for the word dictionary).

\paragraph*{Binary} For a binary part the planner first checks byte-plane widths whose planes are more predictable
than the bytes. For parts of at least 256\,KiB it scores a set of candidates on a 1\,MiB slice from the middle of the
part: x86 code (density of E8/E9 call and jump opcodes), 16-bit images, fixed-size records, float32 and float16
arrays (an exponent byte far more predictable than the mantissa), integer arrays of width 2, 4 and 8, decimal-valued
doubles, and 8-bit images (by a search for the row length). The cheapest candidate wins if it is at least 3\% below
plain coding of the slice. Four further 32\,KiB probes elsewhere in the part detect mixed content, in which case the
choice is verified per block.

\paragraph*{Per-block coding} The plan is file- or part-level, but its application is per block: a block whose
transform does not pay is stored with the plan's transform flag cleared, a block that does not shrink is stored raw,
and at the higher levels several coders are tried per block and the smallest kept. Blocks are 2--16\,MiB at the LZ
levels (about eight balanced blocks per input at the default level) and up to 128\,MiB at the block-sorting levels.

\todo{Add a small figure: the decision flow from part to plan to per-block coder (from docs/ROUTING-MAP.md
\emph{How the router decides}).}

% ======================================================================
\section{Transforms and graph blocks}\label{sec:transforms}

\paragraph*{Typed tables} The table transform (plan 3) splits delimited text into columns and types each value as an
integer, a decimal with a scale, or text. Numeric columns are coded as differences within a group (for example,
consecutive rows of the same entity, found by a column whose value changes rarely) or as residuals against a
\emph{formula} of other columns (a lag, a sum, a ratio with a fitted constant), found by testing candidate formulas
on row subsets and rejecting early those that fail. Irregular rows are kept verbatim with their positions. The
typed transpose (plan 1) is a cheaper variant for sparse integer matrices.

\paragraph*{Logs and JSON} The log transform (plan 4) separates each line into a template of constant tokens and a
stream of variable fields, which are typed and coded per field position. The JSON tokenisers treat keys and
punctuation as the template and every scalar as a field; minified JSON is first split into records.

\paragraph*{Executables and images} The x86 filter converts relative call and jump targets to absolute ones, in the
manner of the BCJ filters of XZ Utils; the ARM64 filter does the same for BL and ADRP and carries the page phase of
each block in the plan, since ADRP arithmetic depends on the address modulo 4096 (without the phase the output was
0.35\% larger\src{docs/METADATA.md}). Image planes are predicted by the median edge detector of
LOCO-I~\todoin{cite LOCO-I / JPEG-LS}.

\subsection{Graph blocks}\label{sec:graph}
For structures that no single transform covers, a block may hold a \emph{graph program} (codec 12). Stream 0 is the
block's bytes. Each node consumes earlier streams and appends its outputs to the stream table; every stream that no
node consumes is a leaf and is coded by an ordinary codec. The decoder decodes the leaves, then runs the nodes'
inverses from last to first. Version 12 defines six node types:
\begin{itemize}
\item \textsc{Fields}: split delimited text into a shape stream and $k$ column streams;
\item \textsc{Decimal}: split a column into tags, scales and mantissas, optionally as differences from the previous
value;
\item \textsc{Concat}: join several streams;
\item \textsc{Int}: integers of width 1--8, with order-0/1/2 prediction, a common divisor, a bias and optional byte
planes;
\item \textsc{Float}: rotate each float word left by one bit so that sign and exponent share the top byte, and
split the words into byte planes;
\item \textsc{Img8}: LOCO-I prediction over rows of a given length and channel step.
\end{itemize}
Leaves are coded by any block coder or by order-0 rANS~\cite{ans}, Huffman-only coding, or a plane-context rANS coder;
the encoder keeps, per leaf, the smallest of these. The planner uses graphs for tables wider than 256 columns, for
integer and float arrays, and for 8-bit images, when a slice codes at least 3\% smaller than the other candidates.
This is close in spirit to OpenZL's graphs, but the node set is closed: a decoder refuses any node or leaf codec it
does not know, and a new node type is a new identifier under the compatibility promise.

% ======================================================================
\section{Coders and levels}\label{sec:coders}\label{sec:levels}

All coders in version 12 are Kolmo's own code or C ports with attribution; no compression library is linked. Five
families serve one level scale (Table~\ref{tab:levels}):
\begin{itemize}
\item \textbf{KLZ}: byte-aligned LZ77 sequences with no entropy stage (the LZ4 role).
\item \textbf{KLZH2}: an LZ77 parse with three repeat offsets, coded as four streams (literals, literal-length codes,
match-length codes, offset codes) each with canonical length-limited Huffman~\cite{huffman} or rANS~\cite{ans}, chosen per stream, plus raw
extra bits (the zstd role). Levels 0--6 use greedy or lazy parses; levels 7--9 a price-based optimal parse over a
binary-tree match finder. Every level writes the same format.
\item \textbf{KLZMA}: an LZMA-class adaptive binary range coder over an optimal parse, written from the published
descriptions of LZMA and not bit-compatible with it.
\item \textbf{BWT}: suffix sorting by libsais~\cite{libsais,sais}, with libbsc's LZP preprocessing and QLFC
second stage~\cite{bsc}, ported to C.
\item \textbf{CM}: an lpaq-style~\cite{lpaq} integer context-mixing coder (hashed contexts of orders 1--20, word,
sparse and record contexts, a match model, small gated mixers and adaptive probability maps), plus image and FASTQ
variants. Integer arithmetic only.
\end{itemize}

\begin{table}[t]
\caption{Level scale of 0.0.2 (docs/LEVELS.md; src/api/api.c). Levels 1--7 decode at LZ speed; 8--9 use block
sorting; ultra is outside the scale.}\label{tab:levels}
\centering\footnotesize
\begin{tabular}{@{}lll@{}}
\toprule
Level & Engine & Analysis \\ \midrule
1 fastest & KLZ (fast parse) & minimal \\
2 & KLZH2 fast parse & minimal \\
3 light (default) & KLZH2 3 & full \\
4 & KLZH2 5 & full \\
5 balanced & KLZH2 7 & full \\
6 & KLZH2 8, 16\,MiB blocks & full \\
7 compact & KLZH2 9 & full \\
8 & BWT for every block & full \\
9 max & BWT (text); BWT or KLZMA (binary) & full \\
U ultra & context mixing & full \\
\bottomrule
\end{tabular}
\end{table}

The level numbers and their \emph{promises} (each level smaller than the one before; levels 1--7 decode at LZ speed)
are a contract; the engine behind a level may change between releases. Two promises do not hold today
(\S\ref{sec:limits}): level~7 is not always smaller than level~6 on small inputs, and levels~8 and~9 give identical
output on plain text.
\todo{Replace Table~\ref{tab:levels} and the level discussion with the tuned assignment from the weak-levels lane
(new level-1 parse; KLZH2 decode loop) once merged; regenerate all tables.}

% ======================================================================
\section{Determinism}\label{sec:determinism}

For a given input, level and release, Kolmo's output is identical across CPUs (x86-64 and ARM64), compilers, thread
counts and memory budgets. Three rules make this hold. First, every coder and planner decision uses integer
arithmetic or floating point that is evaluated identically everywhere (context mixing is integer-only; the build
disables floating-point contraction, and the lzbench integration turns off \texttt{-ffast-math} for Kolmo's objects).
Second, multi-threaded encoding splits work only in ways whose result does not depend on scheduling: blocks are
independent, and the multi-threaded optimal parse of KLZH2 partitions its match finder by hash bucket. Third,
\emph{memory changes speed, never bytes}: a memory governor adapts the number of blocks and trials in flight to the
memory available, but every format decision---block size, table sizes of context mixing, the codec and transform per
block---depends only on the input, the level and an explicit memory cap.

A determinism test compresses a fixed, seeded set of inputs at every level with one and four threads and under forced
budgets of 64\,MiB, 512\,MiB, 4\,GiB and unlimited, and compares the SHA-256 of every frame with a checked-in list
(152 digests). The release binaries are also bit-for-bit reproducible across checkout paths on the three CI
platforms. Determinism is not a compression result, but it is what makes compressed output usable as a cache or
deduplication key within a release and what makes regressions visible.

% ======================================================================
\section{Evaluation}\label{sec:eval}

\subsection{Setup}\label{sec:setup}
All measurements were taken on 2026-11-23 on an Apple M5 (10 cores, on AC power) under macOS, with Kolmo built from
tag \texttt{v0.0.2} with generic compiler flags. Every tool ran on one thread, and every round trip was checked byte
for byte; there were no mismatches. For the standard corpora we used lzbench~2.4.1~\cite{lzbench}, the in-memory
harness most codec authors use, with Kolmo added through a plugin that calls the library (its sizes match the
command line byte for byte) and all codecs built with lzbench's flags. Fast codecs were measured best of three
(\texttt{-t3,3}), slow ones once. The corpora are Silesia~\cite{silesia,silesia_web} as one 212\,MB tar and file by file,
enwik8~\cite{ltcb} (100\,MB of Wikipedia), and Canterbury~\cite{canterbury,canterbury_web} as one 2.8\,MB tar. Speeds are in MiB/s
and ratios are original over compressed size.

\todo{x86-64 rerun (lane x86-bench, \texttt{bench/lzbench/}): replace every lzbench table and
Fig.~\ref{fig:pareto} with x86-64 numbers, keep the M5 numbers as a second platform. First check: Kolmo -1/-3/-9
sizes on Canterbury are identical on x86-64 (determinism), speeds will differ. Note that the x86 lane uses a neutral
Canterbury tar of 2,826,240\,B; the M5 tar here is 2,846,208\,B.}

\subsection{Standard corpora}

\begin{figure*}[t]
\centering
\input{generated/fig-pareto.tex}
\ref{paretolegend}
\caption{Silesia (212\,MB tar), one thread, Apple M5: compression ratio against compression speed (left) and
decompression speed (right) for 50-odd codec settings in lzbench~2.4.1. Kolmo's levels are labelled; the dashed line
is the Pareto front of all points. Source: \texttt{data/lzb-fast-silesia.tar.csv},
\texttt{data/lzb-slow-silesia.tar.csv}. \todoin{x86-64 rerun}}\label{fig:pareto}
\end{figure*}

\begin{table}[t]
\caption{Silesia, one thread, lzbench 2.4.1, Apple M5. C and D: compression and decompression speed in MiB/s.
Source: \texttt{data/lzb-\{fast,slow\}-silesia.tar.csv}. \todoin{x86-64 rerun; tuned levels}}\label{tab:silesia}
\centering\scriptsize\setlength{\tabcolsep}{3.2pt}
\input{generated/tab-silesia.tex}
\end{table}

Table~\ref{tab:silesia} and Fig.~\ref{fig:pareto} place each Kolmo level next to its nearest rival on Silesia. At
the default level~3 Kolmo reaches the ratio of \texttt{zstd -12} (\numSilKThreeR{} against
\numSilzstdOneTwoR{})\src{lzb-*-silesia} and is \numSilThreeVsZstdPct\% smaller than \texttt{zstd -3}, but it
compresses at \numSilKThreeC{} against \numSilzstdThreeC{}\,MiB/s and decodes at \numSilThreeDecFrac{} of zstd's
speed. Level~4 comes within \numSilFourVsZnineteenPct\% of \texttt{zstd -19}'s ratio at
\numSilFourVsZnineteenEnc{} times its compression speed. Level~6 is slightly smaller than \texttt{xz -9}
(\numSilKSixR{} against \numSilxzNineR{}), compresses \numSilSixVsXzEnc{} times faster and decodes
\numSilSixVsXzDec{} times faster\src{lzb-slow-silesia}. Level~9 matches \texttt{zpaq -5}'s ratio
(\numSilKNineR{} against \numSilzpaqFiveR{}) at \numSilNineVsZpaqEnc{} times its compression speed and
\numSilNineVsZpaqDec{} times its decompression speed. Ultra has the highest ratio of the lineup, \numSilKOneZeroR,
at about \numSilKOneZeroC\,MiB/s each way. The fast tier is weak: level~1 is beaten by \texttt{lzav -1} on ratio
(\numSilKOneR{} against \numSillzavOneR) at the same compression speed.

\begin{table}[t]
\caption{enwik8 and Canterbury, one thread, lzbench 2.4.1, Apple M5. Source:
\texttt{data/lzb-\{fast,slow\}-\{enwik8,canterbury.tar\}.csv}. \todoin{x86-64 rerun}}\label{tab:corpora}
\centering\scriptsize\setlength{\tabcolsep}{3pt}
\input{generated/tab-corpora.tex}
\end{table}

On enwik8 (Table~\ref{tab:corpora}) there is little structure to find. Level~3 is \numEnwKThree{} against
\numEnwZThree{} for \texttt{zstd -3}; levels 8 and 9 choose the same block-sorting route and give identical output
(\numEnwKNineBytes\,B), tying libbsc; ultra reaches \numEnwKUBytes\,B (ratio \numEnwKU)\src{lzb-slow-enwik8}, ahead
of the lzbench lineup but well behind the context-mixing leaders of the Large Text Compression
Benchmark~\cite{ltcb}. On Canterbury's small files the low levels lose: level~3 gives \numCantKThree{} against
\numCantZThree{} for \texttt{zstd -3}, and level~7 (\numCantKSeven) is below level~6
(\numCantKSix)\src{lzb-*-canterbury.tar}. Levels 9 and ultra still lead their rivals.

\subsection{Where the gain comes from}

\begin{table}[t]
\caption{Silesia file by file, ratio, one thread. $\Delta_3$: Kolmo -3 against zstd -3. Source:
\texttt{data/lzb-silesia-files.csv}.}\label{tab:files}
\centering\scriptsize\setlength{\tabcolsep}{2.6pt}
\input{generated/tab-silesia-files.tex}
\end{table}

Table~\ref{tab:files} splits Silesia by file. The gain of level~3 over \texttt{zstd -3} is concentrated where the
planner finds structure: the SAO star catalogue (fixed records, $+\numFilesaoGain$\%), the 12-bit X-ray and MRI
images ($+\numFilexrayGain$\% and $+\numFilemrGain$\%), the chemical database (a delimited table,
$+\numFilenciGain$\%) and the x86 DLL ($+\numFileoofficeGain$\%)\src{lzb-silesia-files}. On text, markup and
databases without a regular layout, level~3 is within 3\% of zstd, and on \texttt{xml} it is 1\% larger. On the
structured files, level~3 is smaller than \texttt{zstd -19} and \texttt{xz -9}.

\subsection{Structured data against OpenZL}\label{sec:openzl}

\begin{table*}[t]
\caption{Structured panel, ratio, one thread, command-line tools, best of three. OpenZL 0.3.0 (\texttt{zli}) ran
with the matching format profile at levels 6 and 9 (the generic \texttt{serial} profile where none matches) and
\emph{without training}; the best of its runs is shown. Kolmo ran with no option but the level. Bold: best of
OpenZL and Kolmo. Lvl: Kolmo's best level. Source: \texttt{data/openzl.json}.
\todoin{trained-OpenZL rematch on public held-out data}}\label{tab:structured}
\centering\scriptsize
\input{generated/tab-structured.tex}
\end{table*}

The second experiment (Table~\ref{tab:structured}) compares Kolmo with OpenZL~0.3.0 on \numPanelFiles{} files:
CSV, integer columns, raw images, NumPy arrays, a pandas pickle, model files (ONNX, safetensors, GGUF), the SAO
catalogue, JSON, logs, a FASTA genome and English text. OpenZL ran with the profile that matches each file's format,
which the user must know; Kolmo ran with no option but the level. Counting results within 1\% as ties, Kolmo's
level~3 wins on \numPanelThreeWin{}, ties on \numPanelThreeTie{} and loses on \numPanelThreeLoss{} files; level~5
wins on \numPanelFiveWin, ties on \numPanelFiveTie{} and loses on \numPanelFiveLoss; Kolmo's best level wins on
\numPanelBestWin{} and ties on \numPanelBestTie, all four ties being model-weight files\src{openzl.json}.
OpenZL's CSV profile scores below \texttt{xz -9} on two of the three CSV files, while Kolmo's table transform with
inter-column formulas reaches 2--3 times OpenZL's ratio there. OpenZL leads level~3 and~5 on the integer column
\texttt{year.i32}, the NumPy image stack and the ONNX model. On model weights every tool, Kolmo included, lands at
about 1.2 (float32) and 1.5 (bfloat16); this is the exponent redundancy that ZipNN and DFloat11 also
exploit~\cite{zipnn,dfloat11}, and Kolmo has no edge there.

This comparison is favourable to Kolmo in two ways that we state plainly. OpenZL ran untrained, and its trained
plans are its intended mode of use. And the panel's files come from data used while developing Kolmo, so some tuning
bias is possible.

\todo{Trained-OpenZL rematch (lane bench-page, \texttt{bench/repro/}): replace Table~\ref{tab:structured} by a panel
on public data new to Kolmo---NOAA GHCN-Daily CSV, GH Archive JSON, NASA HTTP logs (1995), NCEP/NCAR reanalysis
float32 grids, GPT-2 / SmolLM2-135M / Qwen2.5-0.5B weights---split into train and test halves by station, hour,
month or year, with OpenZL trained by \texttt{zli train} on the train half (greedy trainer, 20-minute cap per
profile) and every tool measured on the test half on one thread. Add ZipNN (run) and DFloat11 (its published method
computed on the CPU, labelled as such) on the weights. Keep the untrained table as a second column group. Rewrite the
paragraph above from the new counts.}

\subsection{Threads and memory}

\begin{table}[t]
\caption{Silesia with command-line tools, one thread and all ten cores, best of three, with peak resident memory.
Source: \texttt{data/mt-result.json}. \todoin{x86-64 rerun}}\label{tab:mt}
\centering\scriptsize\setlength{\tabcolsep}{2.4pt}
\input{generated/tab-mt.tex}
\end{table}

Because Kolmo's blocks are independent, decoding scales with cores, which zstd's single-frame decoder does not
(Table~\ref{tab:mt}): on ten cores level~3 decodes at \numMtKolmoThreeDecAll{} against \numMtZstdThreeDecAll\,MiB/s
for \texttt{zstd -3}, \numMtThreeVsZstdDecAll{} times faster\src{mt-result.json}. The price is memory: level~9
compresses Silesia with \numMtKolmoNineEncRss\,MB of resident memory on one thread, against
\numMtXzNineEncRss\,MB for \texttt{xz -9}, and level~7 with \numMtKolmoSevenEncRss\,MB. The command line runs at
somewhat different speeds from lzbench because it includes file I/O and process start-up.

% ======================================================================
\section{Limitations}\label{sec:limits}

\paragraph*{Small inputs} Detection needs data: the binary candidates are scored only on parts of at least 256\,KiB,
wide tables and FASTQ need 64\,KiB, and the text trials need enough lines to see a regular layout. A small input
gets little more than the general coders. On Canterbury, level~3 is
\numCantKThree{} against \numCantZThree{} for \texttt{zstd -3}, and level~2 is \numCantKTwo{} against
\numCantZOne{} for \texttt{zstd -1}. In chunked array stores (HDF5 with 1\,MiB chunks, Blosc2 with its automatic
block size) level~1 sometimes does not compress an array at all \todoin{numbers from the hdf5-blosc lane once its PR
lands}.

\paragraph*{Decode speed} At levels 2--7 Kolmo decodes 30--40\% more slowly than zstd on one thread
(\numSilKThreeD{} against \numSilzstdThreeD\,MiB/s at level~3 on Silesia). Its speed advantage is against LZMA and
block sorting, and on many cores. Kolmo is not a replacement for zstd on latency-sensitive paths.

\paragraph*{Fast tier and level ordering} Level~1 is beaten by lzav~\cite{lzav} on ratio and decode speed. Level~7 is not
smaller than level~6 on Canterbury, the result of an optimal parse that is unstable on heterogeneous small inputs;
levels 8 and 9 coincide on text. These violate the level promises of \S\ref{sec:levels}.

\paragraph*{Memory} The ratio-oriented levels use much more memory than their rivals when compressing (above), and
ultra needs about 26 table bytes per input byte with its text contexts, comparable to \texttt{zpaq -m5}.

\paragraph*{Model weights and plain text} Kolmo has no specific model for either beyond the float graph and the
general coders; its results there are at parity with the general tools.

\paragraph*{Evaluation scope} One machine (an ARM64 laptop); a structured panel that is not independent of
development; OpenZL untrained; ZipNN, DFloat11, cmix, paq8 and NNCP not run. \todoin{each of these is addressed by a
running lane or must stay listed here}.

\paragraph*{Project} Kolmo is maintained by one person, is at version 0.0.2, and has no published production users.
The format promise and the specification reduce the risk of storing data with it, but adoption risk remains.

% ======================================================================
\section{Conclusion}

Kolmo shows that per-block detection and trial-based routing, with no user input, recover a large part of what
format-specific configuration buys on tables, logs, numeric arrays, images and executables, while staying a
competitive general-purpose coder on Silesia: archive-class ratios at LZ-class decode speed at levels 4--7, and
ratios at or above the block-sorting and context-mixing archivers at levels 8--9. It does so inside a container
whose decisions are recorded per block, whose decode memory is declared up front, and whose output is deterministic
across platforms. It does not help on small inputs, plain text at fast levels or model weights, and it decodes more
slowly than zstd on one core. The open question we most want to answer is how much of the gap to a \emph{trained}
format-aware compressor remains, which our next measurements address.

\section*{Availability}
Source, format specification, benchmark scripts and the raw data behind every table:
\url{https://github.com/adnaneben/kolmo} (tag \texttt{v0.0.2}; this paper in \texttt{docs/paper/}).

\bibliographystyle{IEEEtran}
\bibliography{refs}

\end{document}
